\documentclass[10pt,a4paper]{amsart}
\usepackage[margin=19mm]{geometry}
\usepackage{amsmath,amssymb,mathtools}
\usepackage[scr=rsfs,cal=euler]{mathalfa}
\usepackage{xfrac}
\usepackage[unicode,hidelinks,bookmarks=false]{hyperref}
\newcommand{\ie}{i.e.\ }
\newcommand{\cf}{cf.\ }
\newcommand{\resp}{respectively\ }
\newcommand{\Subsection}{\subsection}
\providecommand{\nobreakdash}{\nobreak\hbox{-}}
\setlength{\emergencystretch}{2em}
\multlinegap0pt
\usepackage{bm}
\usepackage{bbm}
\usepackage{dsfont}
\usepackage{xspace}
\usepackage{cancel}
\usepackage{mathtools}
\usepackage{amsthm}
\makeatletter
\@ifundefined{theorem}{\newtheorem{theorem}{Theorem}}{}
\@ifundefined{definition}{\newtheorem{definition}[theorem]{Definition}}{}
\@ifundefined{example}{\newtheorem{example}[theorem]{Example}}{}
\@ifundefined{lemma}{\newtheorem{lemma}[theorem]{Lemma}}{}
\makeatother
\newcommand{\Ha}{\mathcal{H}}
\title[Large-scale behaviour of Sobolev functions in Ahlfors regula]{Large-scale behaviour of Sobolev functions in Ahlfors regular metric measure spaces}
\author{Josh Kline, Pekka Koskela, Khanh Nguyen}
\date{Source: arXiv:2310.11718v3 (CC BY 4.0)}
\begin{document}
\maketitle

\noindent\emph{Source and licence.} Large-scale behaviour of Sobolev functions in Ahlfors regular metric measure spaces. Version arXiv:2310.11718v3 (CC BY 4.0), distributed under the Creative Commons Attribution 4.0 International licence, as recorded in the arXiv licence metadata for this exact version and shown on the abstract page https://arxiv.org/abs/2310.11718v3. Full rights notice: licenses/arxiv-2310.11718-CC-BY-4.0.txt.\par\smallskip

\noindent\textit{Editorial scope.} Counted unit: the Example whose source label is \texttt{example5.4-0709} in the article (source0.tex lines 2742--2747, source label \texttt{example5.4-0709}), delivered together with the article's own complete original proof of it (source0.tex lines 2749--2836, one \texttt{proof} environment of 88 lines). No mathematics is written, repaired or re-derived: statement and proof are the article's own text line for line. Required context retained uncounted: thm2.1-0506 (lines 784--801); lem:Measure-Content Bound (lines 977--982); def:almost-thick (lines 1131--1141). Every definition, display and label of those lines is retained unchanged. Omitted from this package and not redistributed: 1--783, 802--976, 983--1130, 1142--2741, 2748--2748, 2837--3069. Nothing inside a retained excerpt is omitted or altered: every retained body is delivered exactly as the article's lines are written, so no omission receipt of any kind accompanies this package. Citations: the retained excerpts carry no citation command at all, so no bibliography entry of the article is transcribed and no reference is redistributed. Reference adaptation: none was needed and none was made; every reference inside the delivered unit resolves inside it, so no unresolved reference or citation is produced. Printed slips kept literal: no printed word, symbol, hypothesis or number of the counted statement or of its proof was altered. Layout and class: the article's own class is \texttt{amsart} with packages including amssymb, amsfonts, amsthm, mathtools, todonotes, enumitem, hyperref, cite, amsmath, bbm, amsxtra, datetime; the wrapper is the lane's pinned \texttt{amsart} editorial wrapper at 10pt on A4, which restates the theorem-like environments the retained text needs and adds the article's own macro definitions that the retained text needs (\texttt{\textbackslash Ha}), with extra wrapper packages (bm, bbm, dsfont, xspace, cancel, mathtools). The delivered numbering is the unit's own. Assets: no retained span contains a figure environment, a graphics-inclusion command, a PDF-page inclusion, a table or a TikZ picture; no image file is copied into this package, the package declares no asset dependency and no image file is redistributed with it. Licence: version arXiv:2310.11718v3 of this article is distributed under the Creative Commons Attribution 4.0 International licence, as recorded in the arXiv licence metadata for this exact version and shown on the abstract page https://arxiv.org/abs/2310.11718v3. A full rights notice is delivered at licenses/arxiv-2310.11718-CC-BY-4.0.txt. Characters: every retained excerpt is pure ASCII, so the delivered engine, XeLaTeX with its default Latin Modern text font, reports no missing character.
\begin{theorem}\label{thm2.1-0506}
Let  $1<Q<+\infty$ and let $1\le p<+\infty$.
Suppose that $(X,d,\mu)$ is an unbounded metric measure space with metric $d$ and Ahlfors $Q$-regular measure $\mu$ supporting a $p$-Poincar\'e inequality.
We assume that $X$ is complete. 
%Let $\mu$ be  $Q$-Ahlfors regular and supports a $Q$-Poincar\'e inequality where $1<Q<\infty$. Suppose that $X$ is complete. 
Then for $0<2r<R<+\infty$, we have that
\[
	\text{\rm Mod}_p(B(O,r), X\setminus B(O,R)) \approx \text{\rm Cap}_p(B(O,r), X\setminus B(O,R))\simeq
 \begin{cases}
     r^{Q-p},&p<Q;\\
     \left(\log \left(\frac{R}{r}\right)\right)^{1-Q},& p=Q;\\
     R^{Q-p},& p>Q;
 \end{cases}  
\]
for a given $O\in X$.
%{\color{red}X is complete ?}
%{\color{red} Would it be better to state this for condensers $(B(x,r),X\setminus B(x,R))$ for $0<r<2r\le R$?}
\end{theorem}

  \begin{lemma}\label{lem:Measure-Content Bound}
  Let $1<Q<+\infty$ and let $1\le p\le Q$.  Suppose that $(X,d,\mu)$ is Ahlfors $Q$-regular, and for $+\infty>\lambda\ge 1$, we denote $A_{\lambda 2^{j+1}}:=B(O,\lambda 2^{j+2})\setminus B(O,\lambda 2^{j+1})$.  Then for $G\subset A_{\lambda 2^{j+1}}$, we have that 
  \[
  \frac{\mu(G)}{(\lambda 2^{j+1})^p}\lesssim\frac{\Ha^{Q-(p-\alpha)}_{+\infty}(G)}{(\lambda 2^{j+1})^\alpha}\quad\text{for all }0<\alpha<p.
  \]      
  \end{lemma}

\begin{definition}\label{def:almost-thick}
 We say that $F$ is  {\it almost thick} if for every $0<\alpha<Q$, there are a finite sequence $\delta_j>0,$ a number $k\in\mathbb N$, and a constant $C\geq 1$ 
 %a sequence $r_j>0$ with $10r_j\leq r_{j+1}$ for all $j\in\mathbb N$ 
 such that 
 \begin{equation}
 \label{almost-thick}
 \mathcal H^{\alpha}_{+\infty} (F\cap A_{C2^{j+1}}) \geq \delta_j (2^{j+1})^\alpha \text{\rm \ \ for all $j\geq k$.
 % and for all $0<\alpha<Q$.
 }
 \end{equation}
\end{definition}

 \begin{example}
 \label{example5.4-0709} Let $1<Q<+\infty$. Suppose that $(X,d,\mu)$ is a complete, unbounded metric measure space with metric $d$ and Ahlfors $Q$-regular measure $\mu$ supporting a $Q$-Poincar\'e inequality.  Then  there is an almost thick set $F$ which is not thick and a function $f\in \dot N^{1,Q}(X)$ such that
 \begin{equation}\label{eq6.3-0809}
 	\text{\rm  $\lim_{F\ni x\to+\infty}f(x)$ exists but\ \ }\lim_{t\to+\infty} f(\gamma(t)) \text{\rm \ \ does not exist for any $\gamma\in\Gamma^{+\infty}$.}
 \end{equation}
 \end{example}

 \begin{proof}Let $O\in X$ be a fixed point, and let 
 %  Let $r_{j}$ be a sequence satisfying $\lim_{j\to+\infty}r_j=+\infty$, $r_j< r_{j+1}$ for all $j\in\mathbb N$ and 
% \begin{equation}\label{eq4.17-0709}
% \lim_{j\to+\infty}\frac{r_{j+1}}{r_j}=+\infty.
% \end{equation}
%  We set $S_j:=\{ x\in X: d(O,x)=r_j+1\}$.  We have the fact that $\mathcal H^{\alpha}_{+\infty}(S_j\cap B(O,r_{j+1})\setminus B(O,r_j))\leq (r_j+2)^\alpha$ for $\alpha\in (0,Q)$ because $S_j$ can be covered by a ball $B(O, r_j+2)$. Combining this with \eqref{eq4.17-0709}, we have that
% \[
% 	\frac{\mathcal H^{\alpha}_{+\infty}(S_j\cap B(O,r_{j+1})\setminus B(O,r_j))}{r_{j+1}^\alpha}\leq \frac{(r_j+2)^\alpha}{r_{j+1}^\alpha}\to 0 \text{\rm \ \ as $j\to+\infty$}.
% \] 
% For any given subsequence $\{S_{j_k}\}_{k\in\mathbb N}$ where $S_{j_k}$ belongs to $\{S_j, j\in\mathbb N \}$,   it then follows that $\cup_{k\in\mathbb N}S_{j_k}$ is non-thick.
% 
% We now pick a correct subsequence $S_{j_k}$ to construct a non-thick set $F$ and $u\in\dot N^{1,Q}(X)$ satisfying the conclusion. 
% 
% {\color{red}write more!}
$\{r_j\}_{j\in\mathbb N}, \{s_j\}_{j\in\mathbb N}$ be two sequences defined by 
 \[
 	r_1=2, s_1=2^{2}, s_{j}=2^{r_{j}}, r_{j+1}=2^{s_j} \ \ \text{\rm for  $j\in\mathbb N$}.
 \]
 Then 
 \begin{equation}\label{eq6.4-0809}
 	100 r_j<s_{j}<r_{j+1}/100 \ \ \text{\rm and \ \ } \lim_{j\to+\infty} \frac{s_j^\alpha}{r_{j+1}^\alpha}=0 \text{\rm \ \ for all $\alpha\in(0,Q)$}.
 \end{equation}
% We set $S_{j}:=\{x\in X: s_j/2\leq d(O,x)\leq s_j\}$ for $j\in\mathbb N$. We have from the fact \eqref{eq2.8-1910} that there is a constant $C>0$ such that   $\mathcal H^{\alpha}_{+\infty}(S_j)\geq C s_j^\alpha$ for all $\alpha\in (0,Q]$ and for all $j\in\mathbb N$. Then
 %\todo[inline]{reference of this face!} Combining this with \eqref{eq6.4-0809}, we obtain that
% \[
% 	\frac{\mathcal H^{\alpha}_{+\infty}(S_j\bigcap B(O,r_{j+1})\setminus B(O,r_j))}{r_{j+1}^\alpha}\geq C \frac{s_j^\alpha}{r_{j+1}^\alpha}=:\delta_j\to0
% \] 
% as $j\to+\infty$ for all $\alpha\in (0,Q)$. It follows that $F:=\cup_{j\in\mathbb N}S_j$ is almost thick. 

%Moreover, we now show that $F$ is not thick. Suppose that $F$ is thick and hence that there are a constant $\delta>0$, $k\in\mathbb N$ such that for all sequence $r_j$ with $10r_j\leq r_{j+1}$ for all $j\in\mathbb N$, we have that
%\[
%	\frac{\mathcal H^{\alpha}_{+\infty}(F\bigcap B(O,r_{j+1})\setminus B(O,r_j))}{r_{j+1}^\alpha}\geq \delta
%\]
%for all $j\geq k$ and for all $0<\alpha<Q$. However, using the sequence $r_j$ defined from the begining, we have that
%\[
%	\frac{\mathcal H^{\alpha}_{+\infty}(F\bigcap B(O,r_{j+1})\setminus B(O,r_j))}{r_{j+1}^\alpha}= \frac{\mathcal H^{\alpha}_{+\infty}(  S_j)}{r_{j+1}^\alpha}\leq \frac{s_j^\alpha}{r_{j+1}^\alpha}\to0 \text{\rm \ \ as $j\to+\infty$}
%\]
%for all $0<\alpha<Q$ which is a contradiction, and so $F$ is not thick. Here the inequality is obtain since $S_j$ is contained in the ball $B(O,s_j)$ and the limit is obtained by \eqref{eq6.4-0809}.

%We only need to construct the function $u$ satisfying \eqref{eq6.3-0809}.  
 We let 
\[\text{
$L_j:=\#\{i\in\mathbb N: r_j\leq 2^i\leq s_j/2\}$ \ and \  $N_j:=\#\{i\in\mathbb N: s_j\leq 2^i\leq r_{j+1}\}$
}
\]
for $j\in\mathbb N$.
 Then $L_j\geq 2^j$ and $N_j\geq 2^j$. We set 
 \[
 g_{s_{j}, r_j}:=  \sum_{i\in\mathbb N, r_j\leq 2^i\leq s_j/2} \frac{2^{-i}}{L_j} \chi_{B(O,2^{i+1})\setminus B(O,2^{i})} \text{\rm \ \ and \ \ } g_{r_{j+1}, s_j}:= \sum_{i\in\mathbb N, s_j\leq 2^i\leq r_{j+1}} \frac{2^{-i}}{N_j}  \chi_{B(O,2^{i+1})\setminus B(O,2^{i})}
 \]
for $j\in\mathbb N$.
 We have that $g:=\sum_{j\in\mathbb N} (g_{s_{j}, r_j}+  g_{r_{j+1}, s_j})$ belongs to $L_\mu^Q(X)$ because 
 \begin{align*}
 \int_Xg^Qd\mu=& \sum_{j\in\mathbb N}  \sum_{i\in\mathbb N, r_j\leq 2^i\leq s_j/2} \left(\frac{2^{-i}}{L_j}\right)^Q \int_{ B(O,2^{i+1})\setminus B(O,2^{i})}d\mu +  \sum_{j\in\mathbb N}  \sum_{i\in\mathbb N, s_j\leq 2^i\leq r_{j+1}} \left(\frac{2^{-i}}{N_j}\right)^Q \int_{ B(O,2^{i+1})\setminus B(O,2^{i})}d\mu \\
 \approx & \sum_{j\in\mathbb N} \frac{1}{L_j^Q}  \sum_{i\in\mathbb N, r_j\leq 2^i\leq s_j/2}  1 + \sum_{j\in\mathbb N}\frac{1}{N_j^Q}  \sum_{i\in\mathbb N, s_j\leq 2^i\leq r_{j+1}}  1\\
 =& \sum_{j\in\mathbb N} \left(\frac{1}{L_j^{Q-1}}+\frac{1}{N_j^{Q-1}}\right)\leq \sum_{j\in\mathbb N} \left(\frac{1}{(2^{j})^{Q-1}} +\frac{1}{(2^j)^{Q-1}}\right)<+\infty. 
 \end{align*}
 Let $\Gamma_{O,x}$ be the collection of all rectifiable curves $\gamma_{O,x}$ connecting $O$ and $x$ where $x\in X$. We define a function $u$ by setting
 \[	
 u(x)=
 \begin{cases}
 	\min\left\{1, \inf_{\gamma\in\Gamma_{O,x}} \int_{\gamma} g_{s_j,r_j}ds \right\}	 & \text{\rm \ \ if \ \ $x\in B(O,s_j/2)\setminus B(O,r_j)$},\\ 
	1& \text{\rm \ \ if \ \ $x\in B(O,s_j)\setminus B(O,s_j/2)$},\\
	1- \min\left\{1, \inf_{\gamma\in\Gamma_{O,x}} \int_{\gamma} g_{r_{j+1},s_j}ds \right\}	& \text{\rm \ \ if \ \ $x\in B(O,r_{j+1})\setminus B(O,s_j)$},
\end{cases}
 \]
for each $j\in\mathbb N$.
 Then $g$ is an upper gradient of $u$ and hence $u\in\dot N^{1,Q}(X)$.  Moreover, $u(x)=0$ for all $x$ such that $d(O,x)=r_j$.  

 For each $10\leq j\in\mathbb N$, we pick $x_j\in X$ so that $B(x_j,2^{-j})\subset B(O,2^{j+1})\setminus B(O,2^j))$, and so that 
 \begin{equation}\label{eq:r_i}B(x_j, 2^{-j})\cap\left(\bigcup_i\{x:|x|=r_i\}\right)=\varnothing.
 \end{equation}
  Since our space is connected by the Poincar\'e inequality,  then we  can find such an $x_j$   by the definition of the $r_i$.  For each $j\ge 10$, define 
 $$B_j:=B(x_j,2^{-j^{\frac{2}{Q-1}}}\times 2^{-j}),$$
 and let $F:=\bigcup_{j=10}^{+\infty}B_j$. By Lemma~\ref{lem:Measure-Content Bound} and Ahlfors $Q$-regularity, we have that  
\begin{align*}
\frac{\mathcal H^{\alpha}_{+\infty}(F\cap B(O,2^{j+1})\setminus B(O,2^j))}{(2^j)^\alpha}=\frac{\Ha^{\alpha}_{+\infty}(B_j)}{(2^j)^\alpha}\gtrsim\frac{\mu(B_j)}{(2^j)^Q}\simeq\frac{(2^{-j^{\frac{2}{Q-1}}}\times2^{-j})^Q}{(2^j)^Q}=:\delta_j\to 0
\end{align*}
as $j\to+\infty$ for all $0<\alpha<Q$.  Thus, $F$ is almost thick as in Definition~\ref{def:almost-thick}. However, $F$ is not thick since 
\[
\frac{\mathcal H^{\alpha}_{+\infty}(F\cap B(O,2^{j+1})\setminus B(O,2^j))}{(2^j)^\alpha}=\frac{\Ha^{\alpha}_{+\infty}(B_j)}{(2^j)^\alpha}\le \frac{(2^{-j^{\frac{2}{Q-1}}}\times2^{-j})^\alpha}{(2^j)^\alpha}\to 0
\]
as $j\to+\infty$ for all $0<\alpha<Q$.

For each $j\geq 10$, we use Theorem~\ref{thm2.1-0506} to choose a function $0\le\psi_j\le 1$ which is admissible for computing $\text{\rm Cap}_Q(B_j, X\setminus B(x_j,2^{-j}))$ such that $$\int_Xg_{\psi_j}^Qd\mu\lesssim \left( \log \left(\frac{1}{2^{-j^{\frac{2}{Q-1}}}}\right)\right)^{1-Q}\approx j^{-2},$$ where $g_{\psi_j}$ is the minimal $Q$-weak upper gradient of $\psi_j$.  By this estimate and the fact that each $\psi_j$ is supported in $B(x_j,2^{-j})$, it follows that $\sum_{j=10}^{+\infty}\psi_j\in  N^{1,Q}(X)$. 

Set $f:=\min\{1, u+\sum_{j=10}^{+\infty}\psi_j\}$.  Since $\psi_j\equiv 1$ on $B_j$, it follows that $f|_F\equiv 1$, and by \eqref{eq:r_i}, we have that $f|_{\{x:|x|=r_i\}}\equiv 0$ for all $i$.  Therefore \eqref{eq6.3-0809} follows, completing the proof. 
 \end{proof}

\end{document}
